\section{2024-11-22 Grundgesetz}

\startframedtipp
In Klausur Bezug auf Grundgesetz
\stopframedtipp

Wie "Toleranz hört bei Intoleranz auf", hört Pazifismus bei gewlalt auf


\startdefinition[title={Subsidiaritätsprinzip}]

\startbuffer[mp:definition:subsidiaritaet]
\startMPpage

u := 1cm;

z0=(0, 2cm);
z1=(0, -2cm);

z2=(4cm, 2cm);
z3=(4cm, 1.5cm);
z4=(4cm, 1cm);
z5=(4cm, 0.5cm);
z6=(4cm, 0);

z7= (0    , 0);
z8= (1.25cm  , 0);
z9= (1.75cm, 0);
z10=(2.27cm, 0);
z11=(2.75cm, 0);

for a=0 upto 4:

	draw z0{dir (-90 - (5 * a))}..(-1cm - 0.5cm * a, 0)..{dir 45}z1;
	draw z0{dir (-90 + (5 * a))}..(1cm + 0.5cm * a, 0)..{dir 135}z1;

endfor

draw z2--z7;
draw z3--z8;
draw z4--z9;
draw z5--z10;
draw z6--z11;

label.rt("Kommune", z2);
label.rt("Land", z3);
label.rt("Bund", z4);
label.rt("EU", z5);
label.rt("UNO", z6);

\stopMPpage
\stopbuffer

\placefigure[force][mp:definition:subsidiaritaet]{Solidiarität subsidiär}{
\typesetbuffer[mp:definition:subsidiaritaet]
}
\stopdefinition

\startframedtipp
Exportierte Waffen dürfen nur unter Detuscher Kontroll und zustimmung genutzt.
\stopframedtipp
